\documentclass[11pt]{article}
\usepackage[margin=1in]{geometry}
\usepackage[T1]{fontenc}
\usepackage[utf8]{inputenc}
\usepackage{lmodern}
\usepackage{hyperref}
\usepackage{booktabs}
\usepackage{longtable}
\usepackage{array}
\usepackage{enumitem}
\usepackage{xcolor}

\hypersetup{
  colorlinks=true,
  linkcolor=blue,
  urlcolor=blue,
  pdftitle={DeepFake Defender Repository Study},
  pdfauthor={Project Technical Study}
}

\setlist[itemize]{leftmargin=1.3em}
\setlist[enumerate]{leftmargin=1.3em}

\title{DeepFake Defender (AI Media Detection Tool)\\Detailed Repository Study: From Initial State to Final Working System}
\author{Technical Project Record}
\date{\today}

\begin{document}
\maketitle
\tableofcontents
\newpage

\section{Purpose of This Document}
This document records the full project journey from the early problem stage to the final working state. It explains:
\begin{enumerate}
  \item what the repository contains,
  \item what problems were found,
  \item what was changed,
  \item why each change was made,
  \item how the final system should be run and maintained.
\end{enumerate}

The writing is intentionally simple and direct, with no heavy wording.

\section{Project Summary}
DeepFake Defender is a multi-layer media analysis system for images and videos. It combines:
\begin{enumerate}
  \item signed provenance (C2PA),
  \item watermark detection (SynthID),
  \item visual model analysis (ResNet + ViT + forensic modules).
\end{enumerate}

The core design principle is to avoid one single detector. Higher-trust evidence is checked first, then lower layers run only when needed.

\section{Initial State and Early Problems}
At the beginning of the work, the system had important behavior issues:
\begin{enumerate}
  \item Real images were sometimes being marked as AI too aggressively.
  \item SynthID-positive files could still continue to model layer, causing conflicting output.
  \item C2PA appeared inconsistent because runtime process and active Python interpreter were not always the same.
  \item UI text and forensic text could disagree with backend decisions.
  \item React home page still showed video detection as ``Coming Soon''.
\end{enumerate}

These were not only quality issues. They also reduced trust in final output.

\section{Project Goals and Constraints}
Main goals during this work:
\begin{enumerate}
  \item Fix logical conflicts in the image decision pipeline.
  \item Reduce false positives without removing core functionality.
  \item Make C2PA behavior reliable and observable.
  \item Keep both frontend and backend messages aligned.
  \item Enable the missing React video flow.
  \item Improve deployment and operations documentation for Windows users.
\end{enumerate}

Constraints:
\begin{enumerate}
  \item Keep existing working functionality intact.
  \item Avoid risky redesign of model logic unless needed.
  \item Prefer small, safe, testable changes.
\end{enumerate}

\section{Repository Structure and Roles}
\subsection{Top-Level}
\begin{itemize}
  \item \texttt{backend/}: Flask app, model orchestration, forensic reporting, template UI.
  \item \texttt{frontend/}: React app and route-based UI pages.
  \item \texttt{model\_output/}: model and support artifacts.
  \item \texttt{requirements.txt}: Python dependencies.
  \item \texttt{README.md}, \texttt{deployment\_document.md}, \texttt{repository\_study.tex}: project and operations documentation.
\end{itemize}

\subsection{Key Backend Files}
\begin{longtable}{p{0.30\textwidth} p{0.64\textwidth}}
\toprule
\textbf{File} & \textbf{Role} \\
\midrule
\endhead
\texttt{backend/app.py} & Main orchestration for routes and final decision wiring. \\
\texttt{backend/combine\_model.py} & Model fusion, thresholds, confidence logic, forensic-aware calibration. \\
\texttt{backend/custom\_forensics.py} & Frequency, anatomy, entropy checks and score shaping. \\
\texttt{backend/forensic.py} & Human-readable summary and report generation. \\
\texttt{backend/src/c2pa\_checker.py} & C2PA runtime checks and manifest parsing. \\
\texttt{backend/video\_detect\_standalone.py} & Video detection path plus fallback behavior. \\
\texttt{backend/serve.py} & Waitress production entrypoint. \\
\bottomrule
\end{longtable}

\subsection{Key Frontend Files}
\begin{longtable}{p{0.30\textwidth} p{0.64\textwidth}}
\toprule
\textbf{File} & \textbf{Role} \\
\midrule
\endhead
\texttt{frontend/src/pages/Home.jsx} & Home selection (image and video). \\
\texttt{frontend/src/pages/Dashboard.jsx} & Image scan dashboard. \\
\texttt{frontend/src/pages/Report.jsx} & Image report output. \\
\texttt{frontend/src/pages/VideoDashboard.jsx} & Video scan dashboard (added). \\
\texttt{frontend/src/pages/VideoReport.jsx} & Video report page (added). \\
\texttt{frontend/src/App.jsx} & React route registration. \\
\texttt{frontend/vite.config.js} & Dev proxy setup to backend. \\
\bottomrule
\end{longtable}

\section{Final Decision Architecture}
\subsection{Image Pipeline (Final)}
Final image logic is strict and layered:
\begin{enumerate}
  \item Layer 1: C2PA check.
  \item Layer 2: SynthID check (if Layer 1 is not decisive).
  \item Layer 3: AI model + forensic support (if still not decisive).
\end{enumerate}

Current policy:
\begin{enumerate}
  \item If any C2PA metadata is present, treat as AI-generated and skip lower layers.
  \item If SynthID watermark is detected, treat as AI-generated and skip model layer.
  \item Otherwise, run model layer and forensic checks.
\end{enumerate}

\subsection{Video Pipeline (Final)}
\begin{enumerate}
  \item Try dedicated video checkpoints first.
  \item If missing or not usable, use frame-sampling fallback with image predictor.
\end{enumerate}

This keeps video detection available even when optional video checkpoints are absent.

\section{Detailed Change Timeline (From Scratch to Final)}
\subsection{Phase 1: Full Codebase Analysis}
Actions:
\begin{enumerate}
  \item Read backend and frontend flows end-to-end.
  \item Mapped each layer and route output.
  \item Located mismatch points between backend decisions and frontend text.
\end{enumerate}

Result:
\begin{enumerate}
  \item A clear list of exact conflict locations was prepared before editing.
\end{enumerate}

\subsection{Phase 2: Pipeline Logic Corrections}
Actions:
\begin{enumerate}
  \item Enforced SynthID short-circuit behavior.
  \item Prevented layer-3 messaging when model was skipped.
\end{enumerate}

Result:
\begin{enumerate}
  \item Conflicting ``watermark detected but model still ran'' behavior was removed.
\end{enumerate}

\subsection{Phase 3: False Positive Reduction}
Actions:
\begin{enumerate}
  \item Increased conservative behavior in threshold resolution.
  \item Tightened forensic support requirements.
  \item Added guardrails for weak-evidence model-only AI outcomes.
\end{enumerate}

Result:
\begin{enumerate}
  \item Real-image false positives were reduced while preserving AI detection capability.
\end{enumerate}

\subsection{Phase 4: C2PA Runtime Reliability Fix}
Actions:
\begin{enumerate}
  \item Verified package installation and API compatibility.
  \item Found mismatch between system Python process and virtual environment process.
  \item Restarted backend with explicit virtual environment interpreter.
\end{enumerate}

Result:
\begin{enumerate}
  \item C2PA and entropy checks became consistent in live responses.
\end{enumerate}

\subsection{Phase 5: New C2PA Policy Change}
Requested policy change:
\begin{enumerate}
  \item If signed C2PA metadata is present, treat as AI-generated even when manifest does not explicitly declare AI.
\end{enumerate}

Actions:
\begin{enumerate}
  \item Updated backend decision gate in \texttt{backend/app.py}.
  \item Updated forensic summary text in \texttt{backend/forensic.py}.
  \item Updated detection method labels in report outputs.
  \item Updated frontend report explanations to match backend policy.
\end{enumerate}

Result:
\begin{enumerate}
  \item Backend, forensic text, and UI labels now match one policy.
\end{enumerate}

\subsection{Phase 6: C2PA Diagnostics Hardening}
Actions:
\begin{enumerate}
  \item Improved C2PA import and runtime status handling.
  \item Added \texttt{/api/health} endpoint with Python interpreter and C2PA diagnostics.
\end{enumerate}

Result:
\begin{enumerate}
  \item Runtime mismatch and package issues can now be diagnosed quickly.
\end{enumerate}

\subsection{Phase 7: React Video Flow Completion}
Actions:
\begin{enumerate}
  \item Replaced ``Coming Soon'' video card in React home.
  \item Added React video dashboard and video report pages.
  \item Added routes for these pages in React router.
  \item Fixed frontend development proxy to backend port 7860.
\end{enumerate}

Result:
\begin{enumerate}
  \item Both image and video flows are now functional in React UI.
\end{enumerate}

\subsection{Phase 8: Safe Codebase Organization}
Actions:
\begin{enumerate}
  \item Added \texttt{backend/src/\_\_init\_\_.py} for cleaner package structure.
  \item Added Windows helper scripts under \texttt{scripts/windows/}:
  \begin{itemize}
    \item \texttt{start-backend-dev.ps1}
    \item \texttt{start-backend-prod.ps1}
    \item \texttt{start-frontend-dev.ps1}
    \item \texttt{health-check.ps1}
  \end{itemize}
\end{enumerate}

Result:
\begin{enumerate}
  \item Startup is easier and less error-prone for Windows users.
\end{enumerate}

\subsection{Phase 9: Documentation Hardening}
Actions:
\begin{enumerate}
  \item Rewrote README with clear version, dependency, setup, and troubleshooting guidance.
  \item Rewrote deployment guide as Windows-first and free-hosting focused.
  \item Expanded this LaTeX study to include full project history and final architecture.
\end{enumerate}

Result:
\begin{enumerate}
  \item Operations are easier and safer for future updates.
\end{enumerate}

\section{File-Level Change Record}
\begin{longtable}{p{0.30\textwidth} p{0.64\textwidth}}
\toprule
\textbf{File} & \textbf{Final-purpose change} \\
\midrule
\endhead
\texttt{backend/app.py} & Added health endpoint, C2PA metadata-presence policy, cleaned imports to package style. \\
\texttt{backend/src/c2pa\_checker.py} & Added runtime diagnostics and robust parsing/error handling. \\
\texttt{backend/forensic.py} & Updated C2PA wording and method labeling to align with new policy. \\
\texttt{backend/templates/report.html} & Updated detection method text for C2PA policy consistency. \\
\texttt{frontend/src/pages/Report.jsx} & Updated C2PA reason text and status display consistency. \\
\texttt{frontend/src/pages/Home.jsx} & Enabled video card and removed placeholder behavior. \\
\texttt{frontend/src/App.jsx} & Added video routes. \\
\texttt{frontend/src/pages/VideoDashboard.jsx} & Added full video scan page and API call flow. \\
\texttt{frontend/src/pages/VideoReport.jsx} & Added full video report page. \\
\texttt{frontend/vite.config.js} & Updated dev API proxy target to port 7860. \\
\texttt{backend/serve.py} & Added production server entrypoint with Waitress. \\
\texttt{requirements.txt} & Added \texttt{waitress} and \texttt{c2pa-python}. \\
\texttt{README.md} & Rewritten in detailed Windows-first format. \\
\texttt{deployment\_document.md} & Rewritten for free deployment from Windows. \\
\texttt{scripts/windows/*} & Added repeatable command scripts for day-to-day operations. \\
\bottomrule
\end{longtable}

\section{Security and Stability Practices Applied}
\begin{enumerate}
  \item Development and production startup are separated.
  \item Health endpoint gives runtime visibility without exposing private media data.
  \item Temporary uploaded files are cleaned after processing.
  \item Extension-based upload filtering remains in place.
  \item Report text is synchronized with backend decisions to avoid user confusion.
\end{enumerate}

\section{Performance and Reliability Notes}
\begin{enumerate}
  \item Models load once at startup to reduce per-request overhead.
  \item Free cloud hosting can still have cold-start delays.
  \item Video fallback path keeps service working when dedicated checkpoints are missing.
  \item Fewer conflicting decision paths improve user trust and operator debugging speed.
\end{enumerate}

\section{Windows Operations Guide (Short Form)}
\subsection{Development}
\begin{enumerate}
  \item Run \texttt{scripts/windows/start-backend-dev.ps1}.
  \item Run \texttt{scripts/windows/start-frontend-dev.ps1}.
  \item Run \texttt{scripts/windows/health-check.ps1}.
\end{enumerate}

\subsection{Production-like Local Run}
\begin{enumerate}
  \item Run \texttt{scripts/windows/start-backend-prod.ps1}.
  \item Validate \texttt{/api/health}, image scan, and video scan.
\end{enumerate}

\section{Validation Summary}
Final validations included:
\begin{enumerate}
  \item backend syntax checks for modified Python files,
  \item frontend production build success,
  \item live API verification for \texttt{/api/health},
  \item live forensic response verification for updated C2PA policy text,
  \item route-level checks for both image and video flows.
\end{enumerate}

\section{Known Limits and Honest Expectations}
\begin{enumerate}
  \item Free deployment plans may sleep when idle.
  \item Cold starts can slow the first request.
  \item Large model artifacts must be present in runtime environment.
  \item C2PA support can vary by cloud runtime packaging details.
\end{enumerate}

\section{Final Outcome}
The final system now has:
\begin{enumerate}
  \item consistent decision logic across backend, forensic summaries, and UI,
  \item reliable runtime diagnostics,
  \item functional image and video flows in both backend templates and React,
  \item safer and clearer Windows-first operational documentation,
  \item organized helper scripts for daily use.
\end{enumerate}

In short, the project moved from partially conflicting behavior to a coherent, traceable, and deployment-ready state while keeping core functionality intact.

\end{document}
